% !TEX program = xelatex
\documentclass[12pt,a4paper]{article}

% ---------- Fonts & Language ----------
\usepackage{fontspec}
\usepackage{xeCJK}
\setmainfont{Times New Roman}
\setCJKmainfont{SimSun}

% ---------- Page Layout ----------
\usepackage{geometry}
\geometry{left=2.5cm,right=2.5cm,top=2.5cm,bottom=2.5cm}
\usepackage{setspace}
\onehalfspacing

% ---------- Common Packages ----------
\usepackage{graphicx}
\usepackage{booktabs}
\usepackage{amsmath,amssymb}
\usepackage{hyperref}
\usepackage{cite}

% ---------- Title Info ----------
\title{基于神经算子的数字大脑模型}
\author{\textbf{小组成员信息}\\
邓博文（2023K8009910008，dengbowen23@mails.ucas.ac.cn）\\
王子逸（2023K8009991010，wangziyi236@mails.ucas.ac.cn）\\
成员三（学号，邮箱）\\
成员四（学号，邮箱）\\
成员五（学号，邮箱）\\
% 若为5人小组，增加一行
}
\date{\today}


\begin{document}
\maketitle

\begin{abstract}

\end{abstract}

\section{背景}
% wzy

\section{神经算子}
% wzy,神经算子单独剔出来讲。这部分为神经算子的介绍

\section{实验设计与实验方法}
% 数据来源、预处理、模型结构、训练策略、评价指标等
\subsection{连接矩阵计算}
\subsubsection{结构连接矩阵}
% dbw
\subsubsection{有效连接矩阵}
% zzk：如何计算功能连接矩阵。原理和步骤。附上论文引用。
\subsubsection{皮层表面拉普拉斯算子}
% dbw
\subsection{仿真数据生成}
% dbw
\subsection{模型训练}
% wzy
\subsection{fMRI数据输入预测}
% lyh
\subsection{可视化与一致性分析}
% ldq
\section{实验结果}
% ldq & dbw

\section{讨论}
% every one （最后再写）
\subsection{结果的意义}
\subsection{不足}
\subsection{展望}

\section{结论}
% 可选

\section*{致谢}
\subsection{分工}
\begin{itemize}
    \item 邓博文：整体分工与进度把控；数据下载；仿真数据生成。
    \item 王子逸：进行模型设计与训练；设计基础模型mlp，神经算子模型fno，deeponet等并进行对比实验；设计实验报告模板并加入编辑报告。
    \item 
\end{itemize}

\begin{thebibliography}{99}
% 参考文献示例（用 cite 引用）
% \bibitem{fno} Li et al. Fourier Neural Operator for Parametric Partial Differential Equations. ICLR.
% \bibitem{deeponet} Lu et al. DeepONet: Learning nonlinear operators for identifying differential equations based on the universal approximation theorem of operators. arXiv:1910.
\end{thebibliography}

\end{document}
